% ==================================================================
% CHAPTER 1: INTRODUCTION
% ==================================================================
\chapter{Introduction}

\section{Problem Statement}
% Trình bày về khó khăn của văn bản Hán Cổ (không dấu, viết liền).
% Nêu sự cần thiết của bài toán Segmentation & Punctuation cho NLP.

\section{Project Objectives}
% Mục tiêu: Xây dựng pipeline tự động hóa với độ chính xác cao và giữ nguyên văn bản gốc (fidelity).

\section{Scope of Study}
% Dữ liệu: EvaHan 2022.
% Mô hình: Xunzi-Qwen3-8B.

% ==================================================================
% CHAPTER 2: BACKGROUND AND RELATED WORK
% ==================================================================
\chapter{Background and Related Work}

\section{Ancient Chinese Processing Tasks}
% Định nghĩa bài toán Word Segmentation và Punctuation dưới góc độ toán học (Sequence Labeling).

\section{Large Language Models for Classical Chinese}
% Giới thiệu dòng Qwen và đặc biệt là Xunzi-Qwen3-8B (được pre-train trên cổ văn).

\section{Parameter-Efficient Fine-Tuning (PEFT)}
\subsection{Low-Rank Adaptation (LoRA)}
% Giải thích ngắn gọn về LoRA: tại sao chọn nó thay vì Full Fine-tune (tiết kiệm VRAM, hiệu quả).

% ==================================================================
% CHAPTER 3: METHODOLOGY
% (Trọng tâm: Dựa trên file xunzi-finetune.ipynb)
% ==================================================================
\chapter{Methodology}

\section{Dataset Preparation}
\subsection{The EvaHan 2022 Dataset}
% Mô tả nguồn dữ liệu, số lượng mẫu train/test.

\subsection{Preprocessing Pipeline}
% Mô tả các bước: 
% 1. Strip POS tags (loại bỏ nhãn từ loại).
% 2. Normalize whitespace.
% 3. Hard Sliding Window Strategy (cắt đoạn 1024 token để train).

\section{Retrieval-Augmented Strategy (SPEADO)} 
% [Quan trọng] Mô tả kỹ thuật SPEADO từ notebook finetune.
\subsection{MinHash LSH Indexing}
% Giải thích cách dùng MinHash để tìm kiếm nhanh các câu tương đồng trong tập train.
\subsection{Dynamic Few-Shot Prompting}
% Mô tả cách chèn "Reference Example" vào prompt nếu độ tương đồng > threshold.

\section{Model Training Configuration}
% Mô tả cấu hình Unsloth:
% - Base Model: Xunzi-Qwen3-8B
% - LoRA Rank: 32, Alpha: 64
% - Target Modules: Attention + MLP (gate_proj, up_proj...)
% - Loss Function: Completion-only Loss (chỉ tính loss trên phần Assistant trả lời).

% ==================================================================
% CHAPTER 4: INFERENCE SYSTEM DESIGN
% (Trọng tâm: Dựa trên file xunzi-finetune-infer.ipynb - Kỹ thuật)
% ==================================================================
\chapter{Inference System Design}

\section{Long-Text Handling Strategy}
\subsection{Keep-Center Sliding Window}
% Giải thích thuật toán inference:
% - Chia văn bản dài thành các chunk.
% - Overlap vùng biên.
% - Chỉ lấy kết quả ở vùng giữa (center) để tránh lỗi biên.

\section{Robustness and Fidelity Mechanism} 
% [Điểm nhấn kỹ thuật] Đây là phần "ăn điểm" về tính kỹ thuật.
\subsection{Problem of Hallucination in LLMs}
% Nêu vấn đề: LLM hay tự ý sửa chữ hoặc thêm bớt ký tự.
\subsection{Post-hoc Robust Alignment Algorithm}
% Mô tả thuật toán dùng `difflib` trong class InferencePipeline:
% - So sánh chuỗi Raw và chuỗi Pred.
% - Ràng buộc cứng: Chỉ chấp nhận thay đổi là dấu câu/space.
% - Khôi phục 100% ký tự gốc nếu model sai.

% ==================================================================
% CHAPTER 5: EXPERIMENTS AND RESULTS
% (So sánh giữa xunzi-baseline-infer và xunzi-finetune-infer)
% ==================================================================
\chapter{Experiments and Evaluation}

\section{Experimental Setup}
\subsection{Baseline Configuration}
% Mô tả cấu hình Zero-shot từ file xunzi-baseline-infer.ipynb (không LoRA, không Retrieval).
\subsection{Evaluation Metrics}
% Định nghĩa các công thức:
% - Seg F1, Punc F1.
% - Joint F1 (Punc + Seg).
% - Fidelity (Character Error Rate) & OSC Score.

\section{Results and Analysis}

\subsection{Performance Comparison}
% Bảng so sánh kết quả:
% | Metric      | Baseline (Zero-shot) | Fine-tuned (Ours) |
% |-------------|----------------------|-------------------|
% | Seg F1      | ...                  | ...               |
% | Punc F1     | ...                  | ...               |
% | Fidelity    | ...                  | ...               |
% -> Nhận xét mức độ cải thiện vượt bậc sau khi fine-tune.

\subsection{Error Analysis}
% Chèn Confusion Matrix (Heatmap) từ file infer.
% Phân tích các cặp dấu câu model hay nhầm (ví dụ: Comma vs Pause).

\subsection{Impact of Input Length}
% Chèn biểu đồ Line chart (F1 score vs Length Bucket) từ file infer.
% Đánh giá độ ổn định của model với câu dài/ngắn.

% ==================================================================
% CHAPTER 6: CONCLUSION
% ==================================================================
\chapter{Conclusion and Future Work}

\section{Summary of Contributions}
% Tổng kết: Đã xây dựng pipeline hoàn chỉnh từ fine-tune đến inference an toàn.

\section{Limitations}
% Ví dụ: Tốc độ suy luận còn chậm do LLM lớn, phụ thuộc chất lượng dữ liệu.

\section{Future Directions}
% Ví dụ: Thử nghiệm Quantization 4-bit sâu hơn, mở rộng sang triều đại khác.

% ==================================================================
% REFERENCES & APPENDIX
% ==================================================================
% \bibliographystyle{plain}
% \bibliography{references}